\section{Cơ sở Lý thuyết}
\label{sec:related_work}

Các bài toán tìm kiếm đường đi và lập kế hoạch chuyển động trên
lưới rời rạc đã thu hút sự chú ý sâu sắc của cộng đồng Trí tuệ
Nhân tạo và Khoa học Máy tính lý thuyết trong nhiều thập kỷ
\cite{russell2020aima, edelkamp2011heuristic}. Trong phần này,
nhóm khảo cứu các nền tảng lý thuyết liên quan trực tiếp đến
đồ án, bao gồm: lý thuyết độ phức tạp của bài toán lăn khối,
các giải thuật tìm kiếm đồ thị kinh điển, lý thuyết thiết kế
hàm đánh giá Heuristic, và các tiêu chuẩn đánh giá thực nghiệm.

\subsection{Lý thuyết Độ phức tạp và Bài toán Lăn khối Đa diện}
\label{subsec:related_complexity}
Khác với các bài toán điều hướng chất điểm truyền thống trong mê cung,
bài toán lăn các vật thể đa diện trên mặt phẳng lưới $2\text{D}$ đặt
ra những ràng buộc chuyển động đặc biệt về hình học và cấu hình không gian.

Buchin và các cộng sự (2007) \cite{buchin2007rolling} đã tiến hành
nghiên cứu sâu sắc về tính chất không gian trạng thái của bài toán lăn
khối lập phương trên lưới ô vuông. Các tác giả chỉ ra rằng khi các mặt
của khối lập phương được gán nhãn hoặc chịu điều kiện tiếp xúc, đồ thị
trạng thái chuyển tiếp có cấu trúc chu trình phức tạp, đòi hỏi số bước
di chuyển lớn hơn đáng kể so với khoảng cách hình học thông thường.
Nghiên cứu tổng quát hơn của Hearn và Demaine (2009) \cite{hearn2009games}
về lý thuyết trò chơi tổ hợp đã hệ thống hóa các mô hình đồ thị chuyển
động có thể đảo ngược và phân loại chúng vào các lớp độ phức tạp tính toán
khác nhau.

Đối với trò chơi Bloxorz, sự kết hợp giữa khối hình chữ nhật
$1 \times 1 \times 2$ (có ba tư thế định hướng khác nhau) và các cơ
chế sàn tương tác (công tắc, cầu lật, gạch mong manh, phân thân) làm
gia tăng đột biến độ phức tạp bài toán. Nghiên cứu mang tính bước ngoặt
của van der Zanden và Bodlaender (2014) \cite{vanderzanden2014bloxorz}
đã chứng minh một cách tường minh rằng: bài toán xác định tính đến
được của Bloxorz tổng quát là \textbf{PSPACE-complete}. Để thiết lập
chứng minh này, các tác giả đã quy giảm từ bài toán thỏa mãn biểu thức
logic định lượng luân phiên (QBF) thông qua việc xây dựng các cổng
logic cơ học (AND, OR, NOT) bằng chính hệ thống công tắc và cầu nối trong
Bloxorz. Kết quả này khẳng định rằng không tồn tại bất kỳ giải thuật
thời gian đa thức nào có thể giải quyết trọn vẹn mọi bản đồ Bloxorz
(trừ khi $\text{P} = \text{PSPACE}$), do đó việc áp dụng các giải thuật
tìm kiếm thông minh và hàm Heuristic cắt tỉa không gian là hoàn toàn tất yếu.

\subsection{Các Giải thuật Tìm kiếm Không Gian Trạng thái Kinh điển}
\label{subsec:related_search}
Để giải quyết bài toán không gian trạng thái được mô hình hóa dưới dạng
đồ thị có hướng hữu hạn $G = (V, E)$, các giải thuật tìm kiếm không
thông tin cung cấp nền tảng chuẩn mực:

\begin{itemize}[--]
    \item \textbf{Tìm kiếm theo Chiều rộng (BFS):}
    Được công bố lần đầu bởi Moore (1959) \cite{moore1959bfs} trong bài
    toán tìm đường thoát mê cung ngắn nhất. BFS duyệt không gian theo từng
    lớp độ sâu thông qua hàng đợi FIFO. Trong đồ thị có trọng số cạnh
    đồng nhất (mỗi bước lật khối tính là 1 bước di chuyển), BFS bảo đảm
    tính đầy đủ và tính tối ưu về số bước đi. Tuy nhiên, nhược điểm chí
    mạng của BFS là độ phức tạp không gian $O(b^d)$, trong đó $b$ là hệ
    số nhánh và $d$ là độ sâu lời giải ngắn nhất, khiến giải thuật dễ cạn
    kiệt bộ nhớ RAM đối với các bản đồ phức tạp \cite{russell2020aima}.

    \item \textbf{Tìm kiếm theo Chiều sâu (DFS):}
    Được hình thức hóa bởi Tarjan (1972) \cite{tarjan1972dfs} trong ngữ
    cảnh xử lý đồ thị tuyến tính. DFS khai thác tối đa một nhánh tìm kiếm
    thông qua cấu trúc ngăn xếp LIFO. Trong các bài toán có không gian
    trạng thái chứa chu trình vô hạn (như việc lăn khối qua lại giữa hai
    ô liền kề), DFS thuần túy sẽ rơi vào vòng lặp vô tận. Do đó, việc
    kết hợp một tập hợp ghi nhận các trạng thái đã duyệt (\texttt{visited})
    là yêu cầu bắt buộc để đảm bảo tính hữu hạn và dừng của giải thuật
    \cite{edelkamp2011heuristic}. Mặc dù tiết kiệm bộ nhớ trong các nhánh
    nông, DFS không bảo đảm tính tối ưu và có xu hướng sinh ra chuỗi hành
    động dài bất thường. Korf (1985) \cite{korf1985ids} đã đề xuất giải
    thuật Tìm kiếm Sâu dần (IDS) để kết hợp ưu thế bộ nhớ tuyến tính của
    DFS với tính tối ưu độ sâu của BFS.

    \item \textbf{Tìm kiếm Chi phí Đồng nhất (UCS):}
    Là một biến thể tương đương với giải thuật đường đi ngắn nhất của
    Dijkstra (1959) \cite{dijkstra1959graph}. Thay vì mở rộng nút theo
    thứ tự độ sâu, UCS sử dụng hàng đợi ưu tiên để luôn mở rộng trạng thái
    có tổng chi phí tích lũy $g(n)$ nhỏ nhất. UCS đạt tính tối ưu toàn cục
    ngay cả khi chi phí di chuyển giữa các ô không bằng nhau. Trong đồ án
    này, UCS đóng vai trò then chốt khi gán trọng số phạt cho các ô gạch
    cam nguy hiểm nhằm mô phỏng tác tử né tránh rủi ro nứt vỡ sàn.
\end{itemize}

\subsection{Tìm kiếm Heuristic và Tính Chất Chấp Nhận Được}
\label{subsec:related_heuristics}
Nhằm khắc phục sự bùng nổ tổ hợp của các phương pháp duyệt mù, giải thuật
A* được giới thiệu bởi Hart, Nilsson và Raphael (1968) \cite{hart1968astar}
đã cách mạng hóa lĩnh vực tìm kiếm heuristic. A* định hướng tìm kiếm bằng
hàm đánh giá tổng thể:
\begin{equation}
    f(n) = g(n) + h(n)
    \label{eq:astar_eval}
\end{equation}
Trong đó, $g(n)$ là chi phí thực tế đã đi từ trạng thái khởi đầu đến nút
$n$, và $h(n)$ là hàm Heuristic ước lượng chi phí tối thiểu từ $n$ đến
trạng thái đích.

Công trình lý thuyết kinh điển của Dechter và Pearl (1985)
\cite{dechter1985astar} cùng chuyên khảo chuyên sâu của Pearl (1984)
\cite{pearl1984heuristics} đã đặt nền móng toán học vững chắc cho tính
tối ưu của A*:
\begin{enumerate}[1.]
    \item \textbf{Tính Chấp nhận được:}
    Hàm Heuristic $h(n)$ được gọi là chấp nhận được nếu nó không bao giờ
    ước tính vượt quá chi phí thực tế thấp nhất $h^*(n)$ để đi từ $n$ đến
    đích, tức là:
    \begin{equation}
        0 \le h(n) \le h^*(n), \quad \forall n \in S
        \label{eq:admissibility}
    \end{equation}
    Khi đồ thị tìm kiếm là dạng cây hoặc giải thuật sử dụng cơ chế mở lại
    nút thích hợp, tính chất chấp nhận được bảo đảm rằng đường đi đầu tiên
    A* tìm đến trạng thái đích luôn là đường đi có tổng chi phí tối ưu
    nhất.

    \item \textbf{Tính Nhất quán:}
    $h(n)$ được gọi là nhất quán nếu với mọi nút $n$ và nút con $n'$ sinh
    ra bởi hành động $a$ có chi phí $c(n, a, n')$, ta luôn có bất đẳng thức
    tam giác:
    \begin{equation}
        h(n) \le c(n, a, n') + h(n')
        \label{eq:consistency}
    \end{equation}
    Tính nhất quán đảm bảo rằng chuỗi giá trị $f(n)$ không giảm dọc theo
    mọi nhánh tìm kiếm, đồng thời nút đầu tiên được lấy ra khỏi hàng đợi
    ưu tiên đã có giá trị $g(n)$ tối ưu cục bộ, loại bỏ nhu cầu phải mở
    lại nút \cite{edelkamp2011heuristic}.
\end{enumerate}

Để thiết kế hàm Heuristic thỏa mãn tính chấp nhận được cho Bloxorz, kỹ thuật
thư giãn bài toán \cite{pearl1984heuristics} thường được áp dụng: bỏ qua các
chướng ngại vật rơi tự do, bỏ qua sự nứt vỡ của gạch cam và giả định khối
có thể lướt hoặc lăn tự do trên mặt phẳng lưới. Trong đồ án này, nhóm
xây dựng hàm Heuristic dựa trên khoảng cách Manhattan chia đôi (vì mỗi bước
lật khối có thể di chuyển trọng tâm $2$ đơn vị) kết hợp phạt sai lệch tư thế
không gian.

\subsection{Phương Pháp Thực Nghiệm và Đánh Giá Hiệu Năng}
\label{subsec:related_benchmarks}
Nhằm đảm bảo tính khách quan và khoa học trong việc so sánh các giải
thuật tìm kiếm, Sturtevant (2012) \cite{sturtevant2012benchmarks} đã đề
xuất bộ tiêu chuẩn đánh giá tìm kiếm trên lưới bao gồm:
\begin{enumerate}[(i)]
    \item Đo lường thời gian thực thi tìm kiếm độc lập với thời gian đồ họa;
    \item Thống kê số nút đã duyệt và số nút đã sinh ra như một thước đo không
    phụ thuộc phần cứng;
    \item Đánh giá mức độ chiếm dụng bộ nhớ đỉnh;
    \item Đánh giá sai số chiều dài đường đi so với lời giải tối ưu.
\end{enumerate}

Tuân thủ phương pháp luận trên, hệ thống của nhóm tách bạch hoàn toàn
bộ giải AI khỏi Ursina Engine \cite{amland2020ursina, panda3d2024manual}.
Việc giám sát bộ nhớ được thực hiện trực tiếp thông qua thư viện chuẩn
\texttt{tracemalloc} của Python \cite{python2024tracemalloc}, và hàng đợi
ưu tiên được tối ưu hóa bằng cấu trúc Min-Heap của thư viện \texttt{heapq}
\cite{python2024heapq}, thiết lập một khung thực nghiệm chuẩn hóa, có thể
tái lập hoàn toàn.
